% Standalone problem document, generated from the adjacent README.md.
% Compile with XeLaTeX. Mathematical content is maintained in Markdown.
\documentclass[11pt,a4paper]{article}
\usepackage[fontsize=10.5pt]{scrextend}
\usepackage[margin=22mm,headheight=15pt,headsep=9mm,footskip=11mm]{geometry}
\usepackage{fontspec}
\setmainfont{texgyrepagella}[Extension=.otf,UprightFont=*-regular,BoldFont=*-bold,ItalicFont=*-italic,BoldItalicFont=*-bolditalic]
\setsansfont{texgyreheros}[Extension=.otf,UprightFont=*-regular,BoldFont=*-bold,ItalicFont=*-italic,BoldItalicFont=*-bolditalic]
\setmonofont{texgyrecursor}[Extension=.otf,UprightFont=*-regular,BoldFont=*-bold,ItalicFont=*-italic,BoldItalicFont=*-bolditalic,Scale=MatchLowercase]
\usepackage{amsmath,amssymb}
\usepackage{unicode-math}
\setmathfont{texgyrepagella-math.otf}
\usepackage{xcolor,microtype,enumitem,needspace,fancyhdr}
\usepackage{bookmark}
\definecolor{ink}{HTML}{17344B}
\definecolor{accent}{HTML}{197D80}
\definecolor{muted}{HTML}{596773}
\hypersetup{colorlinks=true,linkcolor=accent,urlcolor=accent,pdftitle={MI-13 - Nobori's
spectral-middle-factor commutator
conjecture},pdfauthor={Open Problems in Numerical Linear Algebra}}
\urlstyle{same}
\setlength{\parindent}{0pt}
\setlength{\parskip}{5pt plus 1pt minus 1pt}
\linespread{1.04}
\setlength{\emergencystretch}{3em}
\widowpenalty=10000
\clubpenalty=10000
\displaywidowpenalty=10000
\allowdisplaybreaks[1]
\setlist{nosep,leftmargin=1.5em}
\providecommand{\tightlist}{\setlength{\itemsep}{0pt}\setlength{\parskip}{0pt}}
\providecommand{\pandocbounded}[1]{#1}
\pagestyle{fancy}
\fancyhf{}
\fancyhead[L]{\sffamily\footnotesize\color{muted}OPEN PROBLEMS IN NUMERICAL LINEAR ALGEBRA}
\fancyhead[R]{\sffamily\bfseries\footnotesize\color{accent}MI-13}
\fancyfoot[L]{\parbox[b]{0.9\textwidth}{\sffamily\fontsize{8}{10}\selectfont\color{muted}Editorial ratings; a bounded search cannot certify that a problem remains unsolved.\\Literature check: 2026-09-08}}
\fancyfoot[R]{\sffamily\footnotesize\color{muted}\thepage}
\renewcommand{\headrulewidth}{0.3pt}
\renewcommand{\headrule}{\hbox to\headwidth{\color{accent}\leaders\hrule height \headrulewidth\hfill}}
\makeatletter
\renewcommand{\section}{\Needspace{5\baselineskip}\@startsection{section}{1}{0pt}{12pt plus 2pt minus 2pt}{4pt}{\sffamily\bfseries\large\color{ink}}}
\renewcommand{\subsection}{\Needspace{4\baselineskip}\@startsection{subsection}{2}{0pt}{9pt plus 1pt minus 1pt}{3pt}{\sffamily\bfseries\normalsize\color{ink}}}
\makeatother
\setcounter{secnumdepth}{0}
\begin{document}
{\sffamily\small\color{muted}Matrix inequalities and norms\par}
\vspace{3pt}
{\raggedright\hyphenpenalty=10000\exhyphenpenalty=10000\sffamily\bfseries\fontsize{21}{24}\selectfont\color{ink}Nobori's
spectral-middle-factor commutator conjecture\par}
\vspace{7pt}
{\sffamily\small\textbf{Difficulty:} challenging\quad\textbf{Importance:} interesting
to the community\par}
\vspace{4pt}
{\color{accent}\hrule height 0.8pt}
\vspace{6pt}
\textbf{Provenance:} explicit conjecture

\subsection{Problem statement}\label{problem-statement}

Let \(m,n\ge2\) be integers. For every \(A,C\in\mathbb C^{m\times n}\)
and \(B\in\mathbb C^{n\times m}\), is

\[\|ABC-CBA\|_F^2\le
2\|B\|_2^2\bigl(\sigma_1(A)^2+\sigma_2(A)^2\bigr)\|C\|_F^2?\]

Here \(\|X\|_F=(\operatorname{tr}(X^*X))^{1/2}\),
\(\|X\|_2=\sigma_1(X)\), and \(\sigma_1(X)\ge\sigma_2(X)\ge\cdots\) are
the singular values. Zero matrices and deficient ranks are allowed.

\subsection{Relevance}\label{relevance}

The expression measures failure of two rectangular factors to commute
through a middle map. A bound using the middle factor's operator norm
and only two singular values of an outer factor would sharpen estimates
for products of rectangular matrices.

\subsection{References}\label{references}

\begin{enumerate}
\def\labelenumi{\arabic{enumi}.}
\tightlist
\item
  M. Nobori, \emph{A generalization of the Böttcher--Wenzel inequality
  for three rectangular matrices}, Linear Algebra and its Applications
  725 (2025), 135--144. Conjecture 3.1, equation (8), in
  \href{https://arxiv.org/html/2506.17365v2}{arXiv:2506.17365v2};
  \href{https://doi.org/10.1016/j.laa.2025.07.005}{journal}.
\item
  M. Nobori, \emph{A note on some upper bounds for the Frobenius norm of
  the \(q\)-deformed commutator}, arXiv:2608.03897v1 (2026), Theorem 1.1
  and §3. \href{https://arxiv.org/html/2608.03897v1}{Primary
  manuscript}.
\end{enumerate}

\subsection{Status check --- 2026-09-08}\label{status-check-2026-09-08}

The original arXiv record remains v2, dated 3 July 2025. Its proved
Theorem 1.1 places the spectral norm on \(C\) and the Frobenius norm on
\(B\); that is a different assertion. The author's August 2026 paper
treats two-factor expressions \(AB-qBA\), and does not claim this
conjecture. Searches used \textit{Nobori Conjecture 3.1},
\textit{three rectangular matrices conjecture proof}, and
\textit{Nobori commutator spectral conjecture}. No resolution was
located. The original conjecture and later theorem statements were
inspected; the subscription journal body was not independently
retrieved.
\end{document}
